\section{Evaluasi}
\label{sec:evaluasi-bab10}

\subsection{Soal Pilihan Ganda}

\begin{enumerate}
    \item Perbedaan utama \asm{ADC} dari \asm{ADD} adalah \dots
          \begin{enumerate}
              \item \asm{ADC} hanya bekerja pada 8 bit
              \item \asm{ADC} ikut menambahkan \flag{CF}
              \item \asm{ADC} tidak mengubah \eng{flag}
              \item \asm{ADC} menyimpan hasil di memori
          \end{enumerate}

    \item Instruksi yang tidak mengubah \flag{CF} adalah \dots
          \begin{enumerate}
              \item \asm{ADD} \item \asm{SUB}
              \item \asm{INC} \item \asm{ADC}
          \end{enumerate}

    \item Operasi \hex{7F} + \hex{01} pada 8 bit menghasilkan
          \dots
          \begin{enumerate}
              \item \flag{CF} = 1, \flag{OF} = 0
              \item \flag{CF} = 0, \flag{OF} = 1
              \item \flag{CF} = 1, \flag{OF} = 1
              \item \flag{CF} = 0, \flag{OF} = 0
          \end{enumerate}

    \item Hasil \asm{MUL} 8 bit diletakkan pada \dots
          \begin{enumerate}
              \item \reg{AL} \item \reg{AH}
              \item \reg{AX} \item \reg{DX}:\reg{AX}
          \end{enumerate}

    \item Sisa hasil \asm{DIV} 8 bit diletakkan pada \dots
          \begin{enumerate}
              \item \reg{AL} \item \reg{AH}
              \item \reg{AX} \item \reg{DX}
          \end{enumerate}
\end{enumerate}

\subsection{Soal Uraian}

\begin{enumerate}
    \item Jelaskan perbedaan \asm{ADD} dan \asm{ADC}, serta
          \asm{SUB} dan \asm{SBB}. Tuliskan contoh penjumlahan
          32 bit yang memerlukan \asm{ADC}.

    \item Jelaskan perbedaan \flag{CF} dan \flag{OF}. Berikan
          contoh 8 bit yang menunjukkan \flag{CF} = 1 tetapi
          \flag{OF} = 0, dan contoh \flag{OF} = 1 tetapi
          \flag{CF} = 0.

    \item Jelaskan aturan operand \asm{MUL} dan \asm{DIV} untuk
          ukuran 8 bit dan 16 bit. Tuliskan pula syarat aman agar
          tidak terjadi \eng{divide overflow}.

    \item Tulislah program yang menghitung $(25 + 17) \times 4$,
          lalu menyimpan hasilnya pada variabel word. Jelaskan
          register yang dipakai dan letak hasil akhirnya.
\end{enumerate}

\subsection{Kisi-kisi Penilaian}

\begin{table}[H]
    \centering
    \caption{Kisi-kisi evaluasi Bab~\ref{chap:aritmatika}.}
    \label{tab:kisi-bab10}
    \small
    \begin{tabular}{@{}p{1.5cm}p{5.0cm}p{3.0cm}c@{}}
        \toprule
        \textbf{Sub-CPMK} & \textbf{Indikator yang diukur} &
        \textbf{Bentuk soal} & \textbf{Bobot} \\
        \midrule
        3.1 & Menerapkan ADD, ADC, INC, SUB, SBB, DEC &
        Pilihan ganda 1, 2; Uraian 1 & 30\% \\
        \addlinespace
        3.1 & Menjelaskan pengaruh operasi terhadap flag &
        Pilihan ganda 3; Uraian 2 & 25\% \\
        \addlinespace
        3.1 & Menerapkan MUL dan DIV serta overflow &
        Pilihan ganda 4, 5; Uraian 3 & 30\% \\
        \addlinespace
        3.1 & Menulis program aritmatika &
        Uraian 4 & 15\% \\
        \bottomrule
    \end{tabular}
\end{table}

\begin{checklistbox}
\vspace{1.1em}
\textbf{Checklist Pencapaian Kompetensi
Bab~\ref{chap:aritmatika}}

Seorang mahasiswa dinyatakan kompeten pada Sub-CPMK 3.1 apabila:
\begin{itemize}
    \item[\checked] dapat menggunakan ADD, ADC, dan INC;
    \item[\checked] dapat menggunakan SUB, SBB, dan DEC;
    \item[\checked] dapat menjelaskan pengaruh operasi terhadap
          \eng{flag};
    \item[\checked] dapat membedakan CF dan OF;
    \item[\checked] dapat menyusun aritmatika 32 bit;
    \item[\checked] dapat menggunakan MUL sesuai aturan operand;
    \item[\checked] dapat menggunakan DIV dan memeriksa overflow;
    \item[\checked] dapat memverifikasi hasil dengan debugger.
\end{itemize}
\end{checklistbox}
